% META: {"id": "TCOMPL-LOG-ME-002", "chapitre": "TCOMPL-LOGARITHME-HISTORIQUE", "type_objet": "methode", "capacites_codes": ["C2"], "methodes": ["M2"], "mode_creation": "ex_nihilo", "sources_inspiration": [], "status": "needs_review"}
% BEGIN-VERIFY
% from sympy import symbols, log, Eq, solve
% x = symbols('x', positive=True)
% assert solve(Eq(log(x) + log(2*x), log(8)), x) == [2]
% assert solve(Eq(3*x - 1, x + 5), x) == [3]
% END-VERIFY
\begin{fichemethode}{M2}{Résoudre équations et inéquations avec l'équation fonctionnelle de ln}

\quandUtiliser{%
  Utiliser cette méthode lorsque l'équation ou l'inéquation contient des
  logarithmes ($\ln(u) + \ln(v)$, $\ln(u) = \ln(v)$, $\ln(u) \leq \ln(v)$...)
  et qu'on peut la transformer grace a $\ln(ab) = \ln a + \ln b$.
}

\pasApas{%
  \begin{enumerate}
    \item Déterminer l'ensemble de définition : chaque argument de $\ln$ doit
      être STRICTEMENT positif (intersection des conditions).
    \item Regrouper les logarithmes avec l'équation fonctionnelle :
      $\ln u + \ln v = \ln(uv)$, $\ln u - \ln v = \ln\!\dfrac{u}{v}$.
    \item Se ramener a $\ln(A) = \ln(B)$ ou $\ln(A) \leq \ln(B)$.
    \item Supprimer les logarithmes : $\ln$ etant strictement croissante,
      $\ln A = \ln B \iff A = B$ et $\ln A \leq \ln B \iff A \leq B$
      (avec $A, B > 0$).
    \item Résoudre l'équation obtenue, puis NE GARDER que les solutions
      appartenant a l'ensemble de définition.
  \end{enumerate}
}

\exempleRedige{%
  Résoudre $\ln(x) + \ln(2x) = \ln 8$.

  \medskip
  \commentaireMarge{Conditions d'existence d'abord.}%
  Il faut $x > 0$ et $2x > 0$, donc $x > 0$.

  \medskip
  \commentaireMarge{$\ln u + \ln v = \ln(uv)$.}%
  $\ln(x) + \ln(2x) = \ln(2x^2)$, donc l'équation s'ecrit
  $\ln(2x^2) = \ln 8$, soit $2x^2 = 8$, c'est-a-dire $x^2 = 4$.

  \medskip
  Les solutions de $x^2 = 4$ sont $-2$ et $2$ ; seule $x = 2$ vérifie $x > 0$.
  $\mathcal{S} = \{2\}$.
}

\pieges{%
  \begin{itemize}
    \item \textbf{Piège 1.} Oublier les conditions d'existence et garder une
      solution parasite (ici $x = -2$).
    \item \textbf{Piège 2.} Écrire $\ln(a+b) = \ln a + \ln b$ : l'équation
      fonctionnelle porte sur le PRODUIT, pas la somme.
    \item \textbf{Piège 3.} Pour une inéquation, oublier que la croissance
      stricte de $\ln$ CONSERVE le sens : pas de retournement d'inégalité
      (contrairement a une multiplication par un négatif).
  \end{itemize}
}

\verifier{%
  Substituer chaque solution dans l'équation initiale (les deux membres doivent
  être définis et égaux) et contrôler qu'elle respecte toutes les conditions
  d'existence.
}

\sEntrainer{\refExos{M2}}

\end{fichemethode}
